\clearpage\thispagestyle{empty}
\pdfbookmark[0]{Quellen zur Besprechung von Schumachers Bestimmung der Mannheimer Polhöhe}{schumacher1817-sources}
\begin{center}\Large\bfseries Quellen zur Besprechung von Schumachers Bestimmung der Mannheimer Polhöhe\end{center}

Die Übertragung folgt Carl Friedrich Gauß’ Besprechung in den \emph{Werken}, Bd.~VI (Göttingen, 1874), S.~597--599. Der Erstabdruck erschien in den \emph{Göttingischen gelehrten Anzeigen}, Stück~76, 12.~Mai 1817, S.~753--756. Die Bibliothèque nationale de France verzeichnet Schumachers Schrift als Quartband mit 56 Seiten, gedruckt 1816 bei J.~F. Schulz in Kopenhagen.

Die dritte Korrekturformel enthält im Erstabdruck (S.~754) $\frac{1}{23}$, in der Werkausgabe $\frac{1}{28}$. Im Erstabdruck stehen die neun Tabelleneinträge untereinander (S.~755). Die Übertragung folgt der Anordnung der Werkausgabe in zwei Gruppen und behält den Zwischenraum im Eintrag \emph{8 3} für $\chi$ im Großen Bären bei. Im Erstabdruck steht vor \emph{Uebrigens} ein Punkt (S.~756). Die Werkausgabe lässt dort einen Zwischenraum.

KI-Übertragung aus dem historischen Deutsch in modernes Deutsch. KI-erstellter Satz.

\href{https://gdz.sub.uni-goettingen.de/id/PPN236007467}{Werkausgabe}. Erstdruck, S.~\href{https://images.p.sub.uni-goettingen.de/iiif/image/gdz:PPN319721507_1817:00000765/full/1650,/0/default.jpg}{753}, \href{https://images.p.sub.uni-goettingen.de/iiif/image/gdz:PPN319721507_1817:00000766/full/1650,/0/default.jpg}{754}, \href{https://images.p.sub.uni-goettingen.de/iiif/image/gdz:PPN319721507_1817:00000767/full/1650,/0/default.jpg}{755}, \href{https://images.p.sub.uni-goettingen.de/iiif/image/gdz:PPN319721507_1817:00000768/full/1650,/0/default.jpg}{756}. \href{https://catalogue.bnf.fr/ark:/12148/cb31334340p}{Schumachers Schrift, BnF-Katalog}. \href{https://katalog.slub-dresden.de/en/id/0-1132488389}{SLUB-Katalog}.

\clearpage
